\chapter{Cadre prudentiel et réglementaire de la valorisation du Best Estimate}

Le présent chapitre expose les éléments du régime Solvabilité~II nécessaires à
la compréhension de la problématique étudiée. Après avoir rappelé la logique
générale du dispositif et son architecture en trois piliers, il présente la
place du \textit{Best Estimate} dans le bilan prudentiel. Une attention
particulière est ensuite accordée aux différences entre méthodes déterministes
et stochastiques, à la valorisation des options et garanties ainsi qu'aux
exigences de gouvernance et de traçabilité. Enfin, les évolutions issues de la
revue récente de Solvabilité~II sont précisées lorsqu'elles présentent un lien
direct avec le sujet du mémoire.

\section{Solvabilité II : principes, objectifs et architecture}

\subsection{De Solvabilité I à une approche fondée sur les risques}

La directive 2009/138/CE, dite Solvabilité~II, a instauré un cadre prudentiel
harmonisé pour les entreprises d'assurance et de réassurance de l'Union
européenne. Adoptée en 2009 et entrée en application le
1\textsuperscript{er} janvier 2016, elle a remplacé un dispositif dans lequel
les exigences de solvabilité reposaient principalement sur des grandeurs
comptables et des coefficients proportionnels au volume d'activité. Le nouveau
régime adopte une approche économique et fondée sur les risques, dont
l'objectif premier est d'assurer une protection adéquate des preneurs et des
bénéficiaires.

Cette évolution se traduit par trois principes structurants. Premièrement, les
actifs et les passifs sont évalués selon une approche économique, cohérente avec
les conditions de marché. Deuxièmement, le capital requis dépend de la nature
et de l'intensité des risques effectivement supportés par l'organisme.
Troisièmement, les exigences quantitatives sont complétées par des règles de
gouvernance, de gestion des risques, de contrôle interne et de communication
financière.

Le bilan prudentiel diffère donc du bilan établi selon les normes comptables.
Les actifs sont, en principe, valorisés au montant auquel ils pourraient être
échangés dans le cadre d'une transaction conclue dans des conditions normales
de concurrence. Les passifs sont évalués au montant auquel ils pourraient être
transférés ou réglés dans les mêmes conditions. Pour les engagements
d'assurance, cette logique conduit à une valorisation prospective des flux
futurs et non à la seule reprise des provisions comptables historiques.

\subsection{Une architecture en trois piliers}

Le cadre Solvabilité~II est organisé autour de trois piliers complémentaires.

Le \textbf{pilier 1} regroupe les exigences quantitatives. Il définit les règles
de valorisation des actifs et des passifs, le calcul des provisions techniques,
la détermination des fonds propres ainsi que les exigences de capital. Deux
niveaux de capital sont notamment distingués : le \textit{Solvency Capital
Requirement} (SCR) et le \textit{Minimum Capital Requirement} (MCR). Le
\textit{Best Estimate} étudié dans ce mémoire relève directement de ce premier
pilier.

Le \textbf{pilier 2} porte sur le système de gouvernance et la gestion des
risques. Il impose notamment des fonctions clés, un système de contrôle interne,
une fonction de gestion des risques et la réalisation régulière d'une
évaluation interne des risques et de la solvabilité, ou ORSA
(\textit{Own Risk and Solvency Assessment}). Ce pilier exige que les méthodes,
les hypothèses et les limites des modèles utilisés soient comprises et
documentées. La capacité à expliquer les écarts entre deux méthodes de
valorisation s'inscrit donc également dans cette dimension qualitative.

Le \textbf{pilier 3} concerne la communication d'informations à l'autorité de
contrôle et au public. Il comprend notamment le reporting prudentiel, le rapport
régulier au contrôleur et le rapport sur la solvabilité et la situation
financière. La transparence sur les bases, les méthodes et les principales
hypothèses de valorisation des provisions techniques constitue un élément
important de ces publications.

Ces trois piliers ne doivent pas être considérés isolément. Une valorisation
quantitative conforme au pilier~1 doit reposer sur une gouvernance robuste au
titre du pilier~2 et être accompagnée d'une information suffisamment claire et
traçable dans le cadre du pilier~3. Cette articulation justifie l'intérêt d'une
analyse qui ne se limite pas à produire un montant de \textit{Best Estimate},
mais cherche également à comprendre les facteurs à l'origine de ce montant et
de ses variations.

\subsection{Le ratio de solvabilité et la place du Best Estimate}

Les organismes doivent détenir des fonds propres éligibles suffisants pour
couvrir leur SCR. Le ratio de couverture du SCR,
couramment appelé ratio de solvabilité, s'écrit :
\[
    \text{Ratio de solvabilité}
    =
    \frac{\text{Fonds propres éligibles au SCR}}
         {\text{SCR}}.
\]
Un ratio égal à \(100\,\%\) signifie que les fonds propres éligibles sont
exactement égaux au capital requis. Le SCR est calibré afin de correspondre à
une valeur en risque des fonds propres de base avec un niveau de confiance de
\(99{,}5\,\%\) sur un horizon d'un an.

Le \textit{Best Estimate} intervient principalement au numérateur de ce ratio,
par l'intermédiaire du bilan économique. Toutes choses égales par ailleurs, une
hausse du \textit{Best Estimate} augmente les provisions techniques, réduit
l'excédent des actifs sur les passifs et diminue ainsi les fonds propres de
base. Sa valorisation peut également influencer certains modules du SCR et les
mécanismes d'absorption des pertes par les provisions techniques. La maîtrise
de son niveau et de sa sensibilité constitue donc un enjeu de pilotage
prudentiel, même lorsque l'objet de l'étude n'est pas directement la prédiction
du ratio de solvabilité.

\section{La valorisation des provisions techniques}

\subsection{Une valorisation économique des engagements}

L'article~76 de la directive Solvabilité~II prévoit que la valeur des
provisions techniques corresponde au montant actuel qu'une entreprise
d'assurance devrait payer si elle transférait immédiatement ses engagements à
une autre entreprise. Leur calcul doit être cohérent avec les informations
fournies par les marchés financiers, prudent, fiable et objectif.

Lorsque les flux associés aux engagements ne peuvent pas être répliqués de
manière fiable à l'aide d'instruments financiers disposant d'une valeur de
marché observable, les provisions techniques sont calculées comme la somme du
\textit{Best Estimate} et de la marge de risque :
\[
    PT = BE + RM,
\]
où \(PT\) désigne les provisions techniques, \(BE\) le
\textit{Best Estimate} et \(RM\) la marge de risque.

Le \textit{Best Estimate} représente la valeur actualisée moyenne des flux
futurs nécessaires au règlement des engagements. La marge de risque correspond
au coût du capital qu'une entreprise de référence devrait supporter pour
reprendre et honorer ces engagements jusqu'à leur extinction. Le présent
mémoire porte sur le premier de ces deux éléments.

\subsection{Définition réglementaire du Best Estimate}

L'article~77 de la directive définit le \textit{Best Estimate} comme la moyenne
pondérée par leur probabilité des flux de trésorerie futurs, compte tenu de la
valeur temporelle de l'argent et de la courbe des taux sans risque pertinente. De manière schématique, il peut être représenté
par :
\[
    BE =
    \mathbb{E}\left[
        \sum_{t=1}^{T}
        D(0,t)\,CF_t
    \right],
\]
où \(CF_t\) désigne le flux net associé aux engagements à la date \(t\),
\(D(0,t)\) le facteur d'actualisation entre la date de valorisation et la date
\(t\), et \(T\) l'horizon d'extinction du portefeuille.

Cette écriture met en évidence trois dimensions essentielles. Le calcul est
\textbf{prospectif}, puisqu'il repose sur les flux futurs et non uniquement sur
des données passées. Il est \textbf{actualisé}, afin de prendre en compte la
valeur temporelle de l'argent. Enfin, il est \textbf{probabilisé}, car il doit
intégrer la distribution des résultats futurs et pas seulement une trajectoire
centrale lorsque celle-ci ne suffit pas à représenter les risques.

Le calcul doit reposer sur des informations actualisées et crédibles, des
hypothèses réalistes ainsi que des méthodes actuarielles et statistiques
adéquates, applicables et pertinentes. La projection doit prendre en compte
l'ensemble des entrées et sorties de trésorerie nécessaires au règlement des
engagements pendant toute leur durée. Le \textit{Best Estimate} est calculé
brut de réassurance, les montants recouvrables au titre des contrats de
réassurance étant évalués séparément.

\subsection{Flux projetés et hypothèses de calcul}

Le règlement délégué (UE) 2015/35 précise les principes de projection des flux
utilisés pour déterminer le \textit{Best Estimate}. Ceux-ci comprennent notamment les prestations dues
aux assurés et aux bénéficiaires, les primes futures incluses dans la frontière
des contrats, les frais nécessaires à la gestion des engagements et les
paiements liés aux options contractuelles.

Les hypothèses utilisées doivent représenter de manière cohérente les facteurs
susceptibles d'affecter les flux futurs. Dans le cas d'un portefeuille
d'assurance vie, elles portent notamment sur :

\begin{itemize}
    \item la mortalité et la longévité ;
    \item les rachats et les autres comportements des assurés ;
    \item les frais futurs ;
    \item les rendements financiers et l'évolution des actifs ;
    \item les garanties contractuelles ;
    \item la participation future aux bénéfices ;
    \item les actions futures de management.
\end{itemize}

Les actions futures de management peuvent par exemple concerner l'allocation
d'actifs, la réalisation de plus-values latentes ou la politique de
revalorisation des contrats. Elles ne peuvent pas être introduites librement
dans les projections : elles doivent être réalistes, cohérentes avec les
pratiques de l'entreprise et formalisées dans un plan approuvé. De même, les
hypothèses de comportement des assurés doivent tenir compte de la façon dont
ceux-ci peuvent réagir aux conditions économiques, notamment à l'écart entre le
taux servi par le contrat et les rendements proposés par des placements
concurrents.

La frontière des contrats détermine les flux futurs qui appartiennent aux
engagements existants à la date de valorisation. Cette délimitation est
essentielle, car elle conditionne les primes, prestations et frais intégrés dans
le calcul. Enfin, les données doivent être appropriées, complètes et exactes.
Le recours à des approximations est possible lorsque les données sont
insuffisantes, mais il doit être justifié et ne pas résulter de défaillances des
processus internes de collecte ou de validation.

\section{Approches déterministe et stochastique}

\subsection{Deux méthodes au service d'un même objectif prudentiel}

La réglementation impose le calcul d'une meilleure estimation probabilisée ;
elle n'érige pas le \textit{Best Estimate} déterministe et le
\textit{Best Estimate} stochastique en deux mesures prudentielles concurrentes.
Le choix de la méthode doit être adapté à la nature des engagements et aux
risques affectant les flux. Le règlement délégué exige que le calcul soit
réalisé de manière transparente, que la méthode et ses résultats puissent être
revus par un expert qualifié et que les dépendances des flux à l'égard
d'événements futurs soient correctement reflétées.

Une méthode déterministe peut être pertinente lorsque les flux sont peu
sensibles à la dispersion des variables économiques ou lorsque les relations
sont suffisamment linéaires. À l'inverse, une méthode stochastique devient
nécessaire lorsque la valeur des flux dépend de manière importante de
l'ensemble des trajectoires possibles, en particulier en présence d'options,
de garanties ou d'actions de gestion dynamiques.

\subsection{Le Best Estimate déterministe}

Le \textit{Best Estimate} déterministe est obtenu en projetant le portefeuille
selon une trajectoire économique unique, généralement construite à partir
d'hypothèses centrales. Il peut être représenté de manière simplifiée par :
\[
    BE_{\mathrm{det}}
    =
    \sum_{t=1}^{T}
    D^{\mathrm{det}}(0,t)\,
    CF_t^{\mathrm{det}}.
\]

Cette méthode permet une lecture directe de l'évolution du portefeuille et se
prête facilement aux analyses de sensibilité. Son coût de calcul est limité,
puisqu'une seule trajectoire est projetée. Elle fournit ainsi un point de
comparaison utile et peut synthétiser une part importante des caractéristiques
du portefeuille.

Sa principale limite provient de l'utilisation d'un scénario unique. Lorsque
les flux futurs sont une fonction non linéaire d'une variable économique
aléatoire \(X\), la relation suivante est généralement vérifiée :
\[
    \mathbb{E}\!\left[f(X)\right]
    \neq
    f\!\left(\mathbb{E}[X]\right).
\]
La projection du scénario moyen ne restitue alors pas nécessairement la
moyenne des valeurs obtenues dans l'ensemble des scénarios. Cette différence
est au cœur de la valeur temps des options et garanties.

\subsection{Le Best Estimate stochastique}

L'approche stochastique repose sur un ensemble de scénarios économiques
représentant différentes évolutions possibles des variables financières. Pour
chaque scénario, le modèle ALM projette les actifs, les passifs, les décisions
de gestion et les comportements des assurés. Le \textit{Best Estimate}
stochastique est ensuite estimé par la moyenne des valeurs actualisées :
\[
    \widehat{BE}_{\mathrm{sto}}
    =
    \frac{1}{N}
    \sum_{n=1}^{N}
    \left[
        \sum_{t=1}^{T}
        D^{(n)}(0,t)\,
        CF_t^{(n)}
    \right],
\]
où \(N\) désigne le nombre de scénarios économiques.

Cette approche permet de représenter les interactions entre les variables
financières et les mécanismes propres aux contrats. Elle tient notamment compte
du fait qu'une garantie produit un coût dans les scénarios défavorables sans
que l'assureur bénéficie nécessairement de manière symétrique des scénarios
favorables. Elle permet également de modéliser les politiques de participation
aux bénéfices, les réalisations de plus-values, les arbitrages d'allocation et
les rachats dynamiques.

La qualité du résultat dépend toutefois de plusieurs éléments : le calibrage du
générateur de scénarios économiques, le nombre de simulations, les hypothèses
de comportement, les règles de gestion et la stabilité numérique du modèle.
Le calcul est donc plus riche, mais également plus coûteux et plus difficile à
analyser qu'une projection déterministe.

\subsection{Options, garanties et valeur temps}

Le règlement délégué impose de prendre en compte toutes les garanties
financières et options contractuelles incluses dans les contrats, ainsi que les
facteurs susceptibles d'influencer leur exercice par les assurés. Pour un portefeuille d'assurance vie, il peut
notamment s'agir :

\begin{itemize}
    \item d'un taux minimum garanti ;
    \item d'un mécanisme de participation aux bénéfices ;
    \item d'une garantie en capital ;
    \item d'une option de rachat ;
    \item d'une faculté d'arbitrage entre supports ;
    \item d'une option de conversion ou de prorogation.
\end{itemize}

La valeur intrinsèque d'une garantie correspond à son coût dans la situation
économique observée à la date de valorisation. Sa valeur temps traduit le coût
supplémentaire lié à l'incertitude sur les évolutions futures et à la
possibilité que la garantie devienne plus contraignante dans certains
scénarios. Dans le cadre opérationnel retenu dans ce mémoire, la TVOG est
approchée par la différence :
\[
    TVOG =
    BE_{\mathrm{sto}}-BE_{\mathrm{det}}.
\]

Cette définition doit être comprise dans le périmètre et avec les conventions
du modèle utilisé. Elle dépend notamment du scénario déterministe de référence,
des règles de participation aux bénéfices et des actions de management
modélisées. L'écart ne constitue donc pas une constante propre au portefeuille :
il varie avec l'environnement économique, les caractéristiques des contrats et
la politique de gestion de l'assureur.

Les facteurs étudiés dans la base de données reflètent précisément ces
mécanismes. Une courbe de taux basse peut accroître le coût des garanties et
réduire les marges financières disponibles. Une augmentation de la volatilité
peut renforcer la valeur temps d'une option en élargissant la distribution des
résultats futurs. Le niveau des spreads, les plus-values latentes, la duration
du passif, les rachats et le taux minimum garanti modifient également la
capacité de l'assureur à servir ses engagements dans les différents scénarios.
L'analyse de leur effet conjoint nécessite donc des méthodes capables de
représenter des relations non linéaires et des interactions.

\section{Gouvernance, validation et traçabilité des calculs}

\subsection{Exigences de transparence et de revue}

Le règlement délégué prévoit que le \textit{Best Estimate} soit calculé de
manière transparente et que la méthode ainsi que les résultats puissent être
revus par un expert qualifié. Cette exigence
implique que les organismes soient en mesure de documenter les données
utilisées, les hypothèses retenues, les choix méthodologiques, les limites du
modèle et les contrôles réalisés.

La transparence ne suppose pas que le modèle soit simple, mais que son
fonctionnement soit suffisamment compris pour permettre une analyse critique.
Un modèle stochastique peut comporter un grand nombre de règles et
d'interactions ; il doit néanmoins être possible d'identifier les principales
sources de variation des résultats et d'expliquer les écarts significatifs.

\subsection{Rôle de la fonction actuarielle}

La fonction actuarielle coordonne le calcul des provisions techniques et
évalue l'adéquation des méthodes et hypothèses utilisées. Elle doit notamment
apprécier l'incertitude associée aux estimations, s'assurer de l'évaluation
appropriée des options et garanties, comparer et justifier les différences
significatives entre les calculs et réaliser des analyses de sensibilité sur
les principaux risques sous-jacents.

Ces missions donnent une portée réglementaire directe à la compréhension des
écarts. Il ne suffit pas de constater qu'un calcul stochastique produit un
montant différent du calcul déterministe : il convient d'identifier les
hypothèses et mécanismes responsables de cette différence, d'en apprécier la
matérialité et de vérifier la cohérence économique du résultat.

\subsection{Apport d'un modèle explicatif complémentaire}

Un modèle de \textit{machine learning} ne remplace ni le modèle ALM ni les
contrôles actuariels exigés par Solvabilité~II. Il constitue un modèle
complémentaire, construit à partir des sorties du modèle principal. Son intérêt
réside dans sa capacité à synthétiser un grand nombre de résultats et à mettre
en évidence les relations entre les caractéristiques d'une configuration et
l'écart de valorisation observé.

Son utilisation doit toutefois respecter des principes similaires de
gouvernance : séparation des données d'apprentissage et de test, contrôle du
surapprentissage, stabilité des résultats, suivi des erreurs, documentation des
variables et analyse de l'importance des facteurs. Les performances
prédictives ne suffisent pas à elles seules. Dans un contexte prudentiel, le
modèle doit également produire des résultats cohérents avec l'analyse
actuarielle et permettre une interprétation des principaux effets.

La démarche retenue dans ce mémoire répond ainsi à un double objectif. Elle
cherche, d'une part, à prédire l'écart normalisé entre les deux valorisations et,
d'autre part, à fournir une grille de lecture des facteurs qui déterminent cet
écart. Cette seconde dimension est essentielle pour que l'outil contribue à la
traçabilité et à la gouvernance plutôt qu'à la seule automatisation du calcul.

\section{Évolutions issues de la revue de Solvabilité II}

\subsection{Un cadre révisé applicable à compter du 30 janvier 2027}

La directive (UE) 2025/2 a modifié la directive Solvabilité~II, notamment en
matière de proportionnalité, de reporting, de mesures relatives aux garanties
de long terme, de risques de durabilité et de surveillance. Les dispositions révisées et les modifications du
règlement délégué s'appliqueront à compter du 30 janvier 2027. Jusqu'à cette date, les exigences actuellement en
vigueur restent applicables.

Cette distinction temporelle est importante pour le présent mémoire. Les
développements précédents décrivent le cadre applicable à la date de rédaction.
Les règles exposées ci-après constituent une évolution future du régime et ne
doivent pas être présentées comme une possibilité générale déjà ouverte à
l'ensemble des organismes.

\subsection{Introduction d'une évaluation déterministe prudente encadrée}

La revue introduit, pour certaines entreprises de petite taille et non
complexes, la possibilité d'utiliser une évaluation déterministe prudente du
\textit{Best Estimate} pour des engagements d'assurance vie comportant des
options et garanties considérées comme non significatives. Le règlement
délégué (UE) 2026/269 encadre strictement cette possibilité.

Son utilisation suppose notamment que les engagements concernés soient
clairement identifiés, que l'entreprise relève de la catégorie des entreprises
de petite taille et non complexes et que la valeur temps des options et
garanties des engagements visés demeure inférieure à \(5\,\%\) du SCR. La
valorisation déterministe prudente correspond alors à la somme :

\begin{enumerate}
    \item d'une meilleure estimation déterministe des engagements concernés ;
    \item d'une majoration stochastique appliquée au SCR de l'entreprise.
\end{enumerate}

La majoration stochastique est fixée par défaut à \(5\,\%\), sauf si
l'entreprise démontre à l'autorité de contrôle qu'un autre pourcentage
représente plus correctement son profil de risque.

Cette nouvelle mesure ne signifie pas que l'approche stochastique perd son rôle
de reference pour la valorisation des options et garanties matérielles. Elle
constitue une mesure de proportionnalité ciblée, réservée à un périmètre
restreint et accompagnée d'un ajustement prudent. Elle confirme néanmoins
l'importance opérationnelle de la relation entre valorisations déterministe et
stochastique. La capacité à mesurer la valeur temps, à identifier les facteurs
qui la rendent significative et à justifier le niveau d'un ajustement devient
d'autant plus pertinente.

La démarche du présent mémoire ne se confond pas avec cette méthode
réglementaire simplifiée. Elle ne vise pas à calculer la majoration prévue par
le règlement délégué, ni à démontrer l'éligibilité d'une entreprise à ce
dispositif. Elle cherche à modéliser, sur le périmètre étudié, l'écart réellement
observé entre les sorties déterministes et stochastiques de SALLTO. Les
évolutions réglementaires renforcent toutefois l'intérêt du sujet en mettant
explicitement en avant la question de la matérialité de la valeur temps et de
la justification d'une approximation déterministe.


Ainsi, Solvabilité~II impose une valorisation prospective, cohérente avec le marché et
fondée sur une moyenne probabilisée des flux futurs. Pour des contrats
d'assurance vie comportant des options, des garanties, une participation aux
bénéfices ou des comportements dynamiques, l'approche stochastique permet de
représenter des effets que la seule projection d'un scénario central ne capture
pas nécessairement. L'écart entre les deux méthodes traduit alors l'effet de
l'incertitude, des non-linéarités et des interactions entre l'actif, le passif
et la politique de gestion.

Les exigences de transparence, de revue et de gouvernance conduisent les
organismes à ne pas seulement produire un montant, mais également à en
comprendre les déterminants. C'est dans cette perspective que s'inscrit le
recours au \textit{machine learning} dans ce mémoire. Le chapitre suivant
présente le modèle ALM SALLTO ainsi que la méthodologie utilisée pour construire
la base de données destinée à l'analyse de l'écart entre
\(BE_{\mathrm{det}}\) et \(BE_{\mathrm{sto}}\).
